% concatenated from BFY6.tex
\documentclass[12pt]{amsart}

\linespread{1.25}

%\usepackage{showlabels}

\usepackage{amsmath,amssymb,mathrsfs}

\usepackage{amsthm} 

\usepackage{srcltx} 
\usepackage{times}
\usepackage{amsmath,amssymb,mathrsfs}


\usepackage{amsthm} 
\usepackage{srcltx} 


 \usepackage[dvips]{graphicx,color}

\usepackage{amsopn,cite,mathrsfs}
\usepackage{amsfonts}
\usepackage{eucal}

\numberwithin{equation}{section}


\newtheorem{thm}{Theorem}[section]
\newtheorem{lemma}[thm]{Lemma}
\newtheorem{cor}[thm]{Corollary}
\newtheorem{prop}[thm]{Proposition}
\newtheorem{defi}[thm]{Definition}

\newtheorem*{thma}{Theorem A}



\newtheorem*{mthm}{Main Theorem}
\newtheorem*{proof-mthm}{Proof of Main Theorem}

\newcommand{\C}{{\mathbb C}}
\newcommand{\Cast}{{{\mathbb C}^\ast}}
\newcommand{\Castinfty}{{{{\mathbb C}}^\infty}}
\newcommand{\A}{{\mathcal A}(\S)}
\newcommand{\Azero}{{\mathcal A}_0(\S)}
\newcommand{\R}{{\mathbb R}}
\newcommand{\cn}{{\mathbb C}^n}
\newcommand{\T}{{\mathbb T}}
\newcommand{\V}{{\mathbb V}}
\newcommand{\p}{{\mathbb P}}
\newcommand{\D}{{\mathbb D}}
\newcommand{\calD}{{\mathcal D}}
\newcommand{\ind}{\int_\D}
\newcommand{\inc}{\int_{\C}}
\newcommand{\inr}{\int_{\R}}
\newcommand{\dla}{d\lambda_\alpha}
\newcommand{\e}{\varepsilon}                       
\newcommand{\cB}{\EuScript B}
\newcommand{\FF}{\mathcal F}
\newcommand{\kam}{K_{\alpha m}}
\newcommand{\re}{{\rm Re}}
 
 \newcommand{\blem}{\begin{lemma}}
 \newcommand{\elem}{\end{lemma}}
\newcommand{\bpve}{\begin{proof}}
\newcommand{\epve}{\end{proof}}

\newtheorem*{thmA}{Theorem}
\newtheorem*{thmB}{Theorem B}



\def\R{\mathbb{R}}
\def\B{\mathbb{B}}
\def\C{\mathbb{C}}
\def\D{\mathbb{D}}
\def\cD{\overline{\mathbb{D}}}
\def\N{\mathbb{N}}
\def\Z{\mathbb{Z}}
\def\S{\mathbb{S}}
\def\Sde{\mathbb{S} \times \mathbb{S} }
\def\K{\mathbb{K}}
\def\cP{\mathcal{P}}

\def\sp{\vspace*{0.2cm}\\}
\def\img{{\text{\rm Im}\,\gamma}}
\def\Log{{\text{\rm \,Log\,}}}
\def\Res{{\text{\rm \,Res\,}}}
\def\cotan{{\text{\rm \,cotan\,}}}
\def\cosec{{\text{\rm \,cosec\,}}}


\usepackage{amssymb}
\usepackage{amsmath}
\usepackage{amsfonts}
\usepackage[colorlinks=true,linkcolor=red]{hyperref}
\setcounter{MaxMatrixCols}{11}
\newtheorem{theorem}{Theorem}
\theoremstyle{plain}
\newtheorem{corollary}{Corollary}
\newtheorem{definition}{Definition}
\newtheorem{example}{Example}
\newtheorem{proposition}{Proposition}
\newtheorem{remark}{Remark}
\numberwithin{equation}{section}



\begin{document}

\title{On the Commutativity of the Berezin Transform}


\author{Alexander Borichev}

\address{Aix-Marseille University, CNRS, I2M, Marseille,
France}

\email{alexander.borichev@math.cnrs.fr}

\author{G\'erard Fantolini}

\address{Aix-Marseille University, CNRS, I2M, Marseille,
France.}

\email{fantolinigerard@gmail.com}

\author{El-Hassan Youssfi}

\address{Aix-Marseille University, CNRS, I2M, Marseille,
France.}

\email{el-hassan.youssfi@univ-amu.fr}




 
%\subjclass{Primary ....}
\keywords{Reproducing kernel; Fock space; Stieltjes moments.}
%\thanks{This paper is in final form and no version of it will be submitted
%for publication elsewhere.}

\begin{abstract}
We consider the commutativity problem for the Berezin transform on weighted Fock spaces. 
Given a real number $m>0$, for every $\alpha >0$ we denote by $B_{\alpha}$ the Berezin transform associated  to the measure  $\mu_{m}^{\alpha}$ 
with density proportional to $e^{-\alpha |z|^m}$ with respect to Lebesgue measure on the complex plane and normalized so that $\mu_{\phi}^{\alpha}(\mathbb C) =1$. 
We show that the commutativity relation $B_{\alpha}B_{\beta}f=B_{\beta}B_{\alpha}f$ holds for all $f\in L^{\infty}(\mathbb C)$ and $\alpha,\beta> 0$  if and only if $m=2$.  \end{abstract}

\maketitle


\section{Introduction and statement of the main result}

We fix a real number $m>0$ and for every $\alpha >0$ we consider the (normalized) measure on the complex plane  
$$
d\mu_{\alpha,m}(z)= \frac{m\alpha^{2/m}}{2\pi\Gamma(2/m)}e^{-\alpha |z|^m}\,dA(z),
$$ 
where  $dA$ is Euclidean area measure on the complex plane $\mathbb C$ and $\Gamma$ is the Euler gamma function.
We denote by $L_{\alpha,m}^{2}(\mathbb C)$ the space of all square integrable functions with respect to the measure $\mu_{\alpha,m}$ and let 
$$
F_{\alpha,m}^{2}(\mathbb C)=L_{\alpha,m}^{2}(\mathbb C)\cap {\rm Hol}(\mathbb C)
$$
denote the corresponding weighted Fock space. The space $F_{\alpha,m}^{2}(\mathbb C)$ is a Hilbert space of entire functions, and we denote by $K_{\alpha,m}(\cdot,\cdot)$ its reproducing kernel,
$$
f(z)=\langle f,K_{\alpha,m}(\cdot,z) \rangle,\qquad f\in F_{\alpha,m}^{2}(\mathbb C),\, z\in \mathbb C.
$$ 
Let $P_+$ be the orthogonal projection from $L_{\alpha,m}^{2}(\mathbb C)$ onto $F_{\alpha,m}^{2}(\mathbb C)$. Given $f\in L^{\infty}(\mathbb C)$, we define the Toeplitz operator $T_f$ on $F_{\alpha,m}^{2}(\mathbb C)$ by the formula 
$$
T_fg=P_+(fg),\qquad g\in F_{\alpha,m}^{2}(\mathbb C).
$$
The Berezin transform of $T_f$ (or the Berezin transform of $f$) is given by
\begin{multline*}
(B_{\alpha,m}T_f)(z)=(B_{\alpha,m}f)(z)=\frac{\langle T_fK_{\alpha,m}(\cdot,z),K_{\alpha,m}(\cdot,z)\rangle}{K_{\alpha,m}(z,z)}
%$$
%The Berezin transform $B_{\alpha,m}$ is defined for all $f\in L^{\infty}(\mathbb C)$ by
%$$
%B_{\alpha,m}f(z)=
\\=\displaystyle\int_{\mathbb C}f(w) \frac{|K_{\alpha,m}(w,z)|^2}{K_{\alpha,m}(z,z)}\,d\mu_{m}^{\alpha}(w), \qquad z\in\mathbb C.
\end{multline*}%$$
For the Berezin transform of operators in a more general context and for the Berezin transform of functions see \cite{B,E1,E2,Z}.

Since $\|K_{\alpha,m}(\cdot,z) \|^2=K_{\alpha,m}(z,z)$, it is clear that $B_{\alpha,m}$ is a contraction on $L^{\infty}(\mathbb C) $ %into itself 
for every $\alpha>0$.

In the case of the classical Fock spaces, that is, when $m=2$, it is known (see, for instance, \cite[Chapter 2]{Z} that the Berezin transform satisfies the commutativity relation $B_{\alpha,2}B_{\beta,2}=B_{\beta,2}B_{\alpha,2}$ on 
$L^{\infty}(\mathbb C)$.
 
%Obviously, the space $L^{\infty}(\mathbb C)$ of the bounded measurable functions on $\mathbb C$ is densely contained in $L_{\alpha, m}^{2}(\mathbb C)$ for all $\alpha>0$.
 
 
The problem we consider here is to find out for which positive real numbers $m$ does the latter commutativity remain true.  The main goal of this paper is to establish the following:

\begin{thmA}\label{main}  Assume that $m>0$  and consider the test functions  $f_\delta(z):=e^{-\delta|z|^m}$, $z\in\mathbb C$, $\delta>0$. Then the following are equivalent
\begin{itemize}
\item[(1)] The commutativity  relation   $B_{\alpha,m}B_{\beta,m}f=B_{\beta,m}B_{\alpha,m}f$  holds for all $f\in L^{\infty}(\mathbb C)$ and  all $\alpha, \beta > 0$.
\item[(2)] The commutativity  relation   $B_{\alpha,m}B_{\beta,m}f_\delta =B_{\beta,m}B_{\alpha,m}f_\delta$  holds for all  $\alpha, \beta, \delta >0.$
\item[(3)] The commutativity  relation   $(B_{\alpha,m}B_{\beta,m}f_\delta)(0) =(B_{\beta,m}B_{\alpha,m}f_\delta)(0)$  holds for all $\alpha, \beta, \delta >0.$
\item[(4)] $m=2.$
\end{itemize}
\end{thmA} 

The proof of this result relies on expansions of the reproducing kernel using the sequence of the Stieltjes even moments $(s_d)$ of the measure $ e^{-\alpha t^{m}}\,dt$ on $[0,+\infty[$ which will be defined in the sequel. For similar arguments involving the Stieltjes moments see \cite{BY,CF,SY}.


\section{ Proof of the  main result}

We start with the equality
$$
\|z^n\|^2_{\alpha,m}=\frac{\Gamma(2(n+1)/m)}{\Gamma(2/m)\alpha^{2n/m}},\qquad n\in \mathbb Z_+.
$$
Therefore, we have 
$$ 
K_{\alpha,m}(z,w)=\sum_{n\ge 0}\frac{(z\overline{w})^n}{\|z^n\|^2_{\alpha,m}}=
\Gamma(2/m)
\sum_{n\ge 0}\frac{(\alpha^{\frac{2}{m}}z\overline{w})^{n}}{\Gamma(2(n+1)/m)}.
$$

Next we define the Stieltjes moments
$$
s_n(\alpha,m):= \alpha^{-2n/m}\Gamma(2(n+1)/m),
$$
and consider an entire function
$$
S_{\alpha,m}(\zeta):= \sum_{n\ge0}\frac{\zeta^n}{s_n(\alpha,m)}. 
$$
Then
$$
K_{\alpha,m}(z,w)=\Gamma(2/m)S_{\alpha,m}(z\overline{w}).
$$

Consequently, the Berezin transform can be expressed as 
\begin{multline*}
(B_{\alpha,m}f)(z)=\int_{\mathbb C}f(w)\frac{|K_{\alpha,m}(w,z)|^2}{K_{\alpha,m}(z,z)}\,d\mu_{\alpha,m}(w)\\
 = \frac{m\alpha^{2/m}}{2\pi S_{\alpha,m}(|z|^2)}  \int_{\mathbb C}f(w) |S_{\alpha,m}(z\overline{w})|^2
 e^{-\alpha |w|^m}\,dA(w).
\end{multline*}

Since 
$$
S_{\alpha,m}(0)=s_0(\alpha,m)^{-1}= \Gamma(2/m)^{-1},
$$
we have
\begin{equation}
(B_{\alpha,m}f)(0)=\frac{m\alpha^{2/m}}{2\pi\Gamma(2/m)}  \int_{\mathbb C}f(w) e^{-\alpha |w|^m}\,dA(w).
\label{e0}
\end{equation}

Furthermore, 
$$
|S_{\alpha,m}(\zeta)|^2=\sum_{k\ge 0}\sum_{0\le d\le k}\frac{\zeta^d\overline{\zeta}^{k-d}}{s_d(\alpha,m)s_{k-d}(\alpha,m)},\qquad \zeta\in\mathbb C.
$$
%In particular, 
%$$
%|S_{\alpha,m}(0)|^2=s_0(\alpha,m)^{-2}= \Gamma(2/m)^{-2}
%$$
%and 
%Next, a
Applying the Berezin transform to the test function $f_{\delta}(w)=e^{-\delta|w|^{m}}$, $\delta\ge 0$, and using pairwise orthogonality of $w^q$, $q\in\mathbb Z_+$, on the circles $\partial D(0,R)$, $R>0$, we obtain 
\begin{multline}
(B_{\alpha,m}f_\delta)(z)\\= 
\frac{m\alpha^{2/m}}{2\pi S_{\alpha,m}(|z|^2)}  \sum_{k\ge 0} 
\sum_{0\le d\le k}\int_{\mathbb C}f_\delta(w)  
\frac{(z\overline{w})^d(w\overline{z})^{k-d}}{s_{d}(\alpha,m)s_{k-d}(\alpha,m)} e^{-\alpha |w|^m}\,dA(w) \\
 =   \frac{m\alpha^{2/m}}{2\pi S_{\alpha,m}(|z|^2)}\sum_{n\ge 0} |z|^{2n}  \int_{\mathbb C}
\frac{ |w|^{2n} }{s_n^2(\alpha,m)}     e^{-(\alpha+\delta) |w|^m}\,dA(w).\label{e1}
\end{multline}


By \eqref{e0} and \eqref{e1} we have 
\begin{multline*}
(B_{\beta,m}B_{\alpha,m}f_\delta)(0) = \frac{m\beta^{2/m}}{2\pi\Gamma(2/m)}   \int_{\mathbb C}  (B_{\alpha,m}f_\delta)(z) e^{-\beta |z|^m}\,dA(z)\\
=  \frac{m^2 (\alpha\beta)^{2/m} }{4\pi^2\Gamma(2/m)}   \sum_{n\ge 0}   \int_{\mathbb C} \frac{|z|^{2n}}{S_{\alpha,m}(|z|^2)}e^{-\beta |z|^m}\,dA(z)\,\int_{\mathbb C}  
\frac{ |w|^{2n}}{s_n^2(\alpha,m)} e^{-(\alpha + \delta)|w|^m}  \,dA(w) \\
= \frac{m (\alpha\beta)^{2/m} }{2\pi\Gamma(2/m)}  \sum_{n\ge 0} \frac1{s_n^2(\alpha,m)}   \left(\frac{1}{\delta+\alpha}\right)^{\frac{2n+2}{m}}\Gamma\left(\frac{2n+2}{m} \right)\int_{\mathbb C}  \frac{  |z|^{2n} e^{-  \beta |z|^m}  }{ S_{\alpha,m}(|z|^2)}  \,  dA(z)\\
 =    \frac{m (\alpha\beta)^{2/m} }{2\pi\Gamma(2/m)} \sum_{n\ge 0}  
  \frac{ \alpha^{4n/m}}{\Gamma\left(\frac{2n+2}{m} \right)}   \left(\frac{1}{\delta+\alpha}\right)^{\frac{2n+2}{m}}\int_{\mathbb C}  \frac{  |z|^{2n} e^{-  \beta |z|^m}  }{ S_{\alpha,m}(|z|^2)}  \,dA(z).
\end{multline*}
(Here we use that $f_\delta$, $B_{\alpha,m}f_\delta$ are positive and bounded so that one can apply the Foubini theorem.)

Thus,
\begin{multline}\label{BBR1}
(B_{\beta,m}B_{\alpha,m}f_\delta)(0) \\=  \frac{m (\alpha\beta)^{2/m}}{2\pi\Gamma(2/m)}\sum_{n\ge 0}  
   \left(\frac{1}{\delta+\alpha}\right)^{(2n+2)/m} \frac{ \alpha^{4n/m}}{\Gamma(\frac{2n+2}{m})} \int_{\mathbb C}  \frac{  |z|^{2n} e^{-  \beta |z|^m}  }{ S_{\alpha,m}(|z|^2)}\,dA(z).
\end{multline}

Set
$$
U_{\alpha,\beta}(n)=\frac{ \alpha^{4n/m}}{\Gamma(\frac{2n+2}{m})} \int_{\mathbb C}  \frac{  |z|^{2n} e^{-  \beta |z|^m}  }{ S_{\alpha,m}(|z|^2)}\,dA(z).
$$
Since the series in \eqref{BBR1} converges for $\delta=0$, we have
\begin{equation}
U_{\alpha,\beta}(n)\le C\alpha^{2n/m},\qquad  n\ge 0.
\label{estar}
\end{equation}


Now we rewrite \eqref{BBR1} as
\begin{equation}
(B_{\beta,m}B_{\alpha,m}f_\delta)(0) \\=  \frac{m (\alpha\beta)^{2/m}}{2\pi\Gamma(2/m)}\sum_{n\ge 0}  
   \left(\frac{1}{\delta+\alpha}\right)^{(2n+2)/m} U_{\alpha,\beta}(n).
\label{e7}
\end{equation}

\begin{lemma}\label{l1} Let $m,\alpha,\beta>0$. If the commutativity relation 
$$ 
(B_{\beta,m}B_{\alpha,m}f_\delta)(0)=(B_{\alpha,m} B_{\beta,m}f_\delta)(0)
$$
holds for every $\delta>0$, then 
\begin{equation}
U_{\alpha,\beta}(n)=U_{\beta,\alpha}(n),\qquad 0\le n<\frac{m}2.
\label{e8}
\end{equation}
\label{lem1}
\end{lemma}

\begin{proof} 
Suppose that for every $\delta>0$ we have 
$(B_{\beta,m}B_{\alpha,m}f_\delta)(0)=(B_{\alpha,m} B_{\beta,m}f_\delta)(0)$. Then by \eqref{e7} we obtain 
\begin{equation}
\sum_{n\ge 0} \left(\frac{1}{\delta+\alpha}\right)^{\frac{2n+2}m}U_{\alpha,\beta}(n)
\\=\sum_{n\ge 0} \left(\frac{1}{\delta+\beta}\right)^{\frac{2n+2}m}U_{\beta,\alpha}(n),\quad \delta>0.\label{e3}
\end{equation}
We set $t=\delta^{-2/m}$ and rewrite \eqref{e3} as 
\begin{multline}
\sum_{n\ge 0} \frac{t^{n+1}}{(1+\alpha t^{m/2})^{(2n+2)/m}}U_{\alpha,\beta}(n)\\=\sum_{n\ge 0} \frac{t^{n+1}}{(1+\beta t^{m/2})^{(2n+2)/m}}U_{\beta,\alpha}(n),\qquad t>0.
\label{e4}
\end{multline}
Then, using \eqref{estar}, we get
$$
\sum_{n\ge 0}t^n(U_{\alpha,\beta}(n)-U_{\beta,\alpha}(n))=O(t^{m/2}), \qquad t\to 0,
$$
and hence, 
$$
U_{\alpha,\beta}(n)=U_{\beta,\alpha}(n),\qquad 0\le n<\frac{m}2.
$$
\end{proof}

\begin{lemma}\label{mnoneven} Let $m\not\in 2\mathbb Z_+$. If $\alpha\not=\beta$, then the commutativity relation 
\begin{equation}
(B_{\beta,m}B_{\alpha,m}f_\delta)(0)=(B_{\alpha,m} B_{\beta,m}f_\delta)(0),  \quad \delta>0,
\label{estar2}
\end{equation}
does not hold.
\end{lemma}

\begin{proof} 
By \eqref{e4} and \eqref{e8}, relation \eqref{estar2} implies that 
\begin{multline*}
\Bigl(\frac{1}{(1+\alpha t^{m/2})^{2/m}}-\frac{1}{(1+\beta t^{m/2})^{2/m}}\Bigr)U_{\alpha,\beta}(0)\\=
\sum_{0<n<\frac{m}2} \Bigl(\frac{t^n}{(1+\beta t^{m/2})^{(2n+2)/m}}-\frac{t^n}{(1+\alpha t^{m/2})^{(2n+2)/m}}\Bigr)U_{\alpha,\beta}(n)\\+
\sum_{n>\frac{m}2}  \Bigl(\frac{t^n}{(1+\beta t^{m/2})^{(2n+2)/m}}U_{\beta,\alpha}(n)-\frac{t^n}{(1+\alpha t^{m/2})^{(2n+2)/m}}U_{\alpha,\beta}(n)\Bigr)\\=o(t^{m/2}),\qquad t\to 0,
\end{multline*}
and hence,
$$
(\beta-\alpha)U_{\alpha,\beta}(0)=o(1),\qquad t\to 0,
$$
which is impossible unless $\alpha=\beta$.
\end{proof}

\begin{lemma}\label{moments} Let $m$ be an even positive integer. Suppose that the commutativity relation $(B_{\beta,m}B_{\alpha,m}f_\delta)(0)=(B_{\alpha,m}B_{\beta,m}f_\delta)(0)$
holds for every $\delta>0$ and every $\alpha,\beta>0$. Then $m=2$. 
\end{lemma}

\begin{proof} Let $m=2p$. By \eqref{e4} we have
$$
\sum_{n\ge 0} \frac{t^n}{(1+\alpha t^p)^{(n+1)/p}}U_{\alpha,\beta}(n)=\sum_{n\ge 0} \frac{t^n}{(1+\beta t^p)^{(n+1)/p}}U_{\beta,\alpha}(n),\qquad t>0.
$$

By the binomial formula, 
$$
(1+x)^q=\sum_{s\ge 0}c_{q,s}x^s,\qquad 0\le x<1,
$$
where 
$$
c_{q,0}=1,\quad c_{q,1}=q,\quad c_{q,2}=q(q-1)/2,
$$
we obtain
\begin{multline}
\sum_{n,s\ge 0}t^{n+sp}\alpha^s c_{-(n+1)/p,s}U_{\alpha,\beta}(n)\\=\sum_{n,s\ge 0}t^{n+sp}\beta^s c_{-(n+1)/p,s}U_{\beta,\alpha}(n),\qquad 0<t<1,
\end{multline}
with both series converging absolutely.

Hence, for every $b\in\mathbb Z_+$ we have
$$
\sum_{n+sp=b} \alpha^s c_{-(n+1)/p,s}U_{\alpha,\beta}(n)=\sum_{n+sp=b}\beta^s c_{-(n+1)/p,s}U_{\beta,\alpha}(n).
$$
By Lemma~\ref{l1} we know already that 
\begin{equation}
U_{\alpha,\beta}(0)=U_{\beta,\alpha}(0),\label{tt1}.
\end{equation}

Taking the values $b=p,2p$ we obtain 
\begin{equation}
U_{\alpha,\beta}(p)-\frac{\alpha}p U_{\alpha,\beta}(0)=U_{\beta,\alpha}(p)-\frac{\beta}p U_{\beta,\alpha}(0),\label{tt2}
\end{equation}
and
\begin{multline}
U_{\alpha,\beta}(2p)-\frac{\alpha(p+1)}p U_{\alpha,\beta}(p)+\frac{\alpha^2(p+1)}{2p^2}U_{\alpha,\beta}(0)\\=U_{\beta,\alpha}(2p)-\frac{\beta(p+1)}p U_{\beta,\alpha}(p)+\frac{\beta^2(p+1)}{2p^2}U_{\beta,\alpha}(0).\label{tt3}
\end{multline}

By \eqref{tt1} and \eqref{tt2} we obtain
\begin{equation}
U_{\alpha,\beta}(p)-U_{\beta,\alpha}(p)= \frac{\alpha-\beta}p U_{\alpha,\beta}(0).
\label{tt4}
\end{equation}

By \eqref{tt1}--\eqref{tt3} we obtain
\begin{multline}
U_{\alpha,\beta}(2p)-U_{\beta,\alpha}(2p)\\=\frac{(\alpha-\beta)(p+1)}p U_{\alpha,\beta}(p)-\frac{(\alpha-\beta)^2(p+1)}{2p^2}U_{\alpha,\beta}(0).
\label{tt5}
\end{multline}

Next we use that 
\begin{multline}
\frac{\partial}{\partial \beta}U_{\alpha,\beta}(n)=\frac{\partial}{\partial \beta}\frac{ \alpha^{2n/p}}{\Gamma(\frac{n+1}{p})} \int_{\mathbb C}  \frac{  |z|^{2n} e^{-  \beta |z|^{2p}}  }{ S_{\alpha,2p}(|z|^2)}\,dA(z)\\=
-\frac{ \alpha^{2n/p}}{\Gamma(\frac{n+1}{p})} \int_{\mathbb C}  \frac{  |z|^{2n+2p} e^{-  \beta |z|^{2p}}  }{ S_{\alpha,2p}(|z|^2)}\,dA(z)\\=-\frac{n+1}{\alpha^2p}\cdot\frac{ \alpha^{(2n+2p)/p}}{\Gamma(\frac{n+p+1}{p})} \int_{\mathbb C}  \frac{  |z|^{2n+2p} e^{-  \beta |z|^{2p}}  }{ S_{\alpha,2p}(|z|^2)}\,dA(z)\\=-\frac{n+1}{\alpha^2p}U_{\alpha,\beta}(n+p),\qquad n\ge 0.
\label{tt6}
\end{multline}
By analogy,
\begin{equation}
\frac{\partial}{\partial \alpha}U_{\beta,\alpha}(n)=-\frac{n+1}{\beta^2p}U_{\beta,\alpha}(n+p),\qquad n\ge 0.
\label{tt7}
\end{equation}

As a consequence, we obtain 
\begin{equation}
\left\{
\begin{aligned}
\frac{\partial}{\partial \alpha}U_{\beta,\alpha}(0)&=-\frac{1}{\beta^2p}U_{\beta,\alpha}(p),\\
\frac{\partial}{\partial \beta}U_{\alpha,\beta}(0)&=-\frac{1}{\alpha^2p}U_{\alpha,\beta}(p),\\
\frac{\partial}{\partial \alpha}U_{\beta,\alpha}(p)&=-\frac{p+1}{\beta^2p}U_{\beta,\alpha}(2p),\\ 
\frac{\partial}{\partial \beta}U_{\alpha,\beta}(p)&=-\frac{p+1}{\alpha^2p}U_{\alpha,\beta}(2p).
\end{aligned}
\right.
\label{start}
\end{equation}

Now, let us evaluate 
$$
H=\frac{\partial^2}{\partial \alpha\partial \beta}(\alpha^2\beta^2 U_{\alpha,\beta}(0)).
$$
First, by \eqref{tt1}, \eqref{tt4}, and \eqref{start} we have
\begin{multline}
H=\frac{\partial}{\partial \alpha}\Bigl[ 2\alpha^2\beta U_{\alpha,\beta}(0) -\frac{\beta^2}p U_{\alpha,\beta}(p)\Bigr]\\=
\frac{\partial}{\partial \alpha}\Bigl[ 2\alpha^2\beta U_{\beta,\alpha}(0) -\frac{\beta^2}p U_{\beta,\alpha}(p)-\frac{(\alpha-\beta)\beta^2}{p^2}U_{\beta,\alpha}(0)\Bigr]\\=
4\alpha\beta U_{\beta,\alpha}(0)-\frac{2\alpha^2}{\beta p}U_{\beta,\alpha}(p)+\frac{p+1}{p^2} U_{\beta,\alpha}(2p)-\frac{\beta^2}{p^2}U_{\beta,\alpha}(0)+\frac{\alpha-\beta}{p^3}U_{\beta,\alpha}(p)\label{tttdif}.
\end{multline}
By \eqref{tt1}, $H$ is stable under permutation of $\alpha$ and $\beta$, 
$$
H=\frac{\partial^2}{\partial \beta\partial \alpha}(\beta^2\alpha^2 U_{\beta,\alpha}(0)),
$$
and, hence,
\begin{multline*}
4\alpha\beta U_{\beta,\alpha}(0)-\frac{2\alpha^2}{\beta p}U_{\beta,\alpha}(p)+\frac{p+1}{p^2} U_{\beta,\alpha}(2p)-\frac{\beta^2}{p^2}U_{\beta,\alpha}(0)+\frac{\alpha-\beta}{p^3}U_{\beta,\alpha}(p)\\=
4\alpha\beta U_{\alpha,\beta}(0)-\frac{2\beta^2}{\alpha p}U_{\alpha,\beta}(p)+\frac{p+1}{p^2} U_{\alpha,\beta}(2p)-\frac{\alpha^2}{p^2}U_{\alpha,\beta}(0)+\frac{\beta-\alpha}{p^3}U_{\alpha,\beta}(p).
\end{multline*}

We obtain that
\begin{multline*}
\frac{p+1}{p^2} (U_{\alpha,\beta}(2p)-U_{\beta,\alpha}(2p))+ \frac{\beta^2}{p^2}U_{\beta,\alpha}(0)-\frac{\alpha^2}{p^2}U_{\alpha,\beta}(0)\\
+\frac{2\alpha^2}{\beta p}U_{\beta,\alpha}(p)-\frac{2\beta^2}{\alpha p}U_{\alpha,\beta}(p)
+\frac{\beta-\alpha}{p^3}U_{\alpha,\beta}(p)
+\frac{\beta-\alpha}{p^3}U_{\beta,\alpha}(p)=0,
\end{multline*}
and by \eqref{tt1} and \eqref{tt5} we get
\begin{multline*}
\frac{(\alpha-\beta)(p+1)^2}{p^3} U_{\alpha,\beta}(p)-\frac{(\alpha-\beta)^2(p+1)^2}{2p^4}U_{\alpha,\beta}(0)+ \frac{\beta^2-\alpha^2}{p^2}U_{\alpha,\beta}(0)\\
+\frac{2\alpha^2}{\beta p}U_{\beta,\alpha}(p)-\frac{2\beta^2}{\alpha p}U_{\alpha,\beta}(p)
+\frac{\beta-\alpha}{p^3}(U_{\alpha,\beta}(p)
+U_{\beta,\alpha}(p))=0,
\end{multline*}

By \eqref{tt4} we obtain
\begin{multline*}
\frac{(\alpha-\beta)(p+1)^2}{p^3} U_{\alpha,\beta}(p)-\frac{(\alpha-\beta)^2(p+1)^2}{2p^4}U_{\alpha,\beta}(0)- \frac{\alpha^2-\beta^2}{p^2}U_{\alpha,\beta}(0)\\
-\frac{2\alpha^2(\alpha-\beta)}{\beta p^2}U_{\alpha,\beta}(0)+\Bigl(\frac{2\alpha^2}{\beta p}
-\frac{2\beta^2}{\alpha p}\Bigr)U_{\alpha,\beta}(p)
\\-\frac{2(\alpha-\beta)}{p^3}U_{\alpha,\beta}(p)+\frac{(\alpha-\beta)^2}{p^4}U_{\alpha,\beta}(0)=0,
\end{multline*}
Dividing by $\alpha-\beta$ we conclude that for $\alpha\not=\beta$ we have 
\begin{multline*}
\Bigl(\frac{(p+1)^2}{p^3} +\frac{2(\alpha^2+\alpha\beta+\beta^2)}{\alpha\beta p}-\frac{2}{p^3}\Bigr)U_{\alpha,\beta}(p) 
\\ =\Bigl(\frac{(\alpha-\beta)(p+1)^2}{2p^4}+\frac{\alpha+\beta}{p^2}+\frac{2\alpha^2}{\beta p^2}-\frac{\alpha-\beta}{p^4}\Bigr)U_{\alpha,\beta}(0).
\end{multline*}
Fix $\alpha=1$. Then
\begin{multline*}
\Bigl(\frac{2\beta}{p}+O(1)\Bigr)U_{\alpha,\beta}(p)  =\Bigl(\Bigl(\frac{1}{p^2}+\frac{1}{p^4}-\frac{(p+1)^2}{2p^4}\Bigr)\beta+O(1)\Bigr)U_{\alpha,\beta}(0)\\=
\Bigl(\frac{(p-1)^2}{2p^4}\beta+O(1)\Bigr)U_{\alpha,\beta}(0),\qquad \beta\to\infty.
\end{multline*}

Suppose now that $m>2$, that is, $p>1$. Then for some $C>0$ and for all $\beta\ge \beta_0$ we have 
$$
\int_{\mathbb C}  \frac{  |z|^{2p} e^{-  \beta |z|^{2p}}  }{ S_{1,2p}(|z|^2)}\,dA(z)\ge C \int_{\mathbb C}  \frac{e^{-  \beta |z|^{2p}}  }{ S_{1,2p}(|z|^2)}\,dA(z).
$$
Since
\begin{align*}
\int_{\mathbb C}  \frac{e^{-  \beta |z|^{2p}}  }{ S_{1,2p}(|z|^2)}\,dA(z)&\ge C\beta^{-1/p},\qquad \beta\to\infty,\\
\int_{\mathbb C}  \frac{|z|^{2p} e^{-  \beta |z|^{2p}}  }{ S_{1,2p}(|z|^2)}\,dA(z)&=O(\beta^{-(p+1)/p}),\qquad \beta\to\infty,
\end{align*}
we obtain a contradiction. Thus, $m=2$. 
\end{proof}


 

\begin{proof}[Proof of Theorem A] The implication $(4)\implies(1)$ is contained in the book of Zhu \cite[Section 3.2]{Z}.
The proof of the remaining parts follows from Lemmata \ref{mnoneven} and \ref{moments}.
\end{proof}

\begin{thebibliography}{1}
%\bibitem{B} H. Bommier-Hato, \textit{Hankel operators and Berezin Transform}, Thesis 2008

\bibitem{B} F.~A.~Berezin, \textit{Covariant and contravariant symbols of operators}, Math.\ USSR-Izv.\ {\bf 6} (1972), 1117--1151.

\bibitem{BY} H.~Bommier-Hato, El.~H.~Youssfi, \textit{Hankel operators and the Stieltjes moment problem}, J. Funct.\ Anal.\ {\bf 258} (2010) 978--998.

\bibitem{CF} C.~Castante, J.~F\`abrega, D.~Pascuas, \textit{Hankel bilinear forms on generalized Fock-Sobolev spaces on $\mathbb C^{n}$}, Ann.\ Acad.\ Sci.\ Fenn.\ Math.\ 
{\bf 45} (2020) 841--862.

\bibitem{E2} M.~Engli\v{s}, \textit{Functions invariant under the Berezin transform}, J. Funct.\ Anal.\ {\bf 121} (1994) 233--254.

\bibitem{E1} M.~Engli\v{s}, \textit{Asymptotics of the Berezin transform and quantization on planar domains}, Duke Math.\ J. {\bf 79} (1995) 57--76.

\bibitem{SY} K.~Seip, El.~H.~Youssfi, \textit{Hankel operators on Fock spaces and related Bergman kernel estimates}, J. Geom.\ Anal.\ {\bf 23} (2013) 170--201.

\bibitem{Z} K.~Zhu, Analysis on Fock spaces, Springer, 2012.





\end{thebibliography}



\end{document}



\begin{lemma}\label{interversion}   For all $\alpha >0, $ we have the following estimate
$$ 
\aligned
\sum_{d=0}^{n}   \frac{1}{s_{d}(\alpha)s_{n-d}(\alpha} = \left\{   \aligned  & O\left(\frac{n \alpha^{\frac{2n}{m}} }{\Gamma\left(\frac{n+1}{m}\right)\Gamma(\frac{n+3}{m})}\right), 
\ \ \text{if \  }  n \   \ \text{is odd },  
\\
& O\left(\frac{n \alpha^{\frac{2n}{m}}}{\left(\Gamma\left(\frac{n+2}{m}  \right)\right)^2} \right), 
\ \ \text{if  \ } n \   \ \text{is even}.
\endaligned
\right.
\endaligned
$$

\end{lemma}

In particular,
$$
\aligned
(\displaystyle B_{\alpha}f)(0) & =    \frac{m}{2\pi} \alpha^{\frac{2}{m}}\int_{\mathbb C} 
f(w) 
  e^{-\alpha |w|^m} 
 dA(w).
\endaligned
$$

Applying the latter equalities to the  function $f_{\delta}(w)=e^{-\delta| w|^{m}}$ where $\delta>0$, yields.

$$
\aligned
\displaystyle B_{\alpha}f_\delta(z)  & =   \frac{m}{2\pi} \alpha^{\frac{2}{m}} \sum_{n=0}^{+\infty}  \frac{ |z|^{2n}  }{S_{\alpha}(|z|^2)}\int_{\mathbb C} 
\frac{ |w|^{2n}}{s_{n}^2(\alpha)} 
  e^{-(\alpha +\delta) |w|^m} 
 dA(w)
\endaligned
$$




Setting $\delta = 1/t$ yields
  $$  \aligned  & \sum_{n=0}^{+\infty}  
  \frac{ \alpha^{\frac{4n}{m}}}{\Gamma\left(\frac{2n+2}{m} \right)}   \left(\frac{t}{1+\alpha t}\right)^{\frac{2n+2}{m}}\int_{\mathbb C}  \frac{  |z|^{2n} e^{-  \beta |z|^m}  }{ S_{\alpha}(|z|^2)}   dA(z)  \\
  &  =  \sum_{n=0}^{+\infty}  
  \frac{ \beta^{\frac{4n}{m}}}{\Gamma\left(\frac{2n+2}{m} \right)}   \left(\frac{t}{ 1+\beta t}\right)^{\frac{2n+2}{m}}\int_{\mathbb C}  \frac{  |z|^{2n} e^{-  \alpha |z|^m}  }{ S_{\beta}(|z|^2)}    dA(z)
  \endaligned
  $$ 
  for all  $ \alpha, \beta, t>0$. 
  
  
  Dividing by $t^{\frac{2}{m}}$ and taking the limit as $t$ tends to 0 in the latter equality gives
  
  
  Multiplying both parts by $\delta^{2/m}$ and passing to the limit $\delta\to \infty$ we obtain 
$$
\int_{\mathbb C}  \frac{e^{-  \beta |z|^m}  }{ S_{\alpha,m}(|z|^2)} \,   dA(z)      =  \int_{\mathbb C}  \frac{   e^{-  \alpha |z|^m}  }{ S_{\beta,m}(|z|^2)}\,  dA(z).
$$ 
  By induction we can prove that
  $$  \aligned  \frac{ \alpha^{\frac{4n}{m}}}{\Gamma\left(\frac{2n+2}{m} \right)}  \int_{\mathbb C}  \frac{ |z|^{2n}e^{-  \beta |z|^{m}}  }{ S_{\alpha}(|z|^2)}    dA(z)    &  = \frac{ \beta^{\frac{4n}{m}}}{\Gamma\left(\frac{2n+2}{m} \right)}  
   \int_{\mathbb C}  \frac{ |z|^{2n}  e^{-  \alpha |z|^m}  }{ S_{\beta}(|z|^2)}   dA(z)
  \endaligned
  $$ 
for all nonnegative integers $n<\frac{m}{2}.$ From this observation it 
  follows that 
  $$  \aligned  &  \sum_{n=0}^{\lfloor \frac{m}{2} \rfloor }  
      \left(    
  \left(\frac{t}{1+\alpha t}\right)^{\frac{2n+2}{m}} -  \left(\frac{t}{1+\beta t}\right)^{\frac{2n+2}{m}}  \right) 
  \frac{ \alpha^{\frac{4n}{m}}}{\Gamma\left(\frac{2n+2}{m} \right)}  \int_{\mathbb C}  
  \frac{  |z|^{2n} e^{-  \beta |z|^{m}  }}{ S_{\alpha}(|z|^2)  }dA(z)    \\
    & =  \sum_{n=\lfloor \frac{m}{2} \rfloor +1}^{+\infty}  
  \frac{ \alpha^{\frac{4n}{m}}}{\Gamma\left(\frac{2n+2}{m} \right)} 
    \left(\frac{t}{1+\alpha t}\right)^{\frac{2n+2}{m}}\int_{\mathbb C} 
     \frac{  |z|^{2n} e^{-  \beta |z|^m}  }{ S_{\alpha}(|z|^2)}   dA(z)   \\
  &  -  \sum_{n=\lfloor \frac{m}{2} \rfloor+1}^{+\infty}  
  \frac{ \beta^{\frac{4n}{m}}}{\Gamma\left(\frac{2n+2}{m} \right)}   \left(\frac{t}{ 1+\beta t}\right)^{\frac{2n+2}{m}}\int_{\mathbb C}  \frac{  |z|^{2n} e^{-  \alpha |z|^m}  }{ S_{\beta}(|z|^2)}    dA(z)
  \endaligned
  $$ 
  Dividing by $t^{\frac{2 +m}{m}}$ and taking the limit as $t$ tends to 0 in the latter equality yields
  $$  \aligned   \frac{2}{m} \left( \alpha - \beta \right)  \int_{\mathbb C}  \frac{e^{-  \beta |z|^m}  }{ S_{\alpha}(|z|^2)}    dA(z)   =  0  \endaligned
  $$ 
This leads to  a contradiction when $\alpha \not = \beta$.

We argue as in the proof of Lemma \ref{mnoneven} and use the same notations as before. It is not hard to see that 
 Assuming that the commutativity relation $ \left(B_{\beta}B_{\alpha}f_\delta\right)(0) = \left(B_{\alpha} B_{\beta}f_\delta\right)(0) $   holds for all  $ \alpha, \beta, \delta>0$ implies that
  $$  \aligned  \frac{ \alpha^{\frac{4n}{m}}}{\Gamma\left(\frac{2n+2}{m} \right)}  \int_{\mathbb C}  \frac{ |z|^{2n}e^{-  \beta |z|^{m}}  }{ S_{\alpha}(|z|^2)}    dA(z)    &  = \frac{ \beta^{\frac{4n}{m}}}{\Gamma\left(\frac{2n+2}{m} \right)}  
   \int_{\mathbb C}  \frac{ |z|^{2n}  e^{-  \alpha |z|^m}  }{ S_{\beta}(|z|^2)}   dA(z)
  \endaligned
  $$ 
for all nonnegative integers $n<\frac{m}{2}.$ 




Hence
  follows that 
  $$  \aligned  &  \sum_{n=0}^{\frac{m}{2} - 1}  
      \left(    
  \left(\frac{t}{1+\alpha t}\right)^{\frac{2n+2}{m}} -  \left(\frac{t}{1+\beta t}\right)^{\frac{2n+2}{m}}  \right) 
  \frac{ \alpha^{\frac{4n}{m}}}{\Gamma\left(\frac{2n+2}{m} \right)}  \int_{\mathbb C}  
  \frac{  |z|^{2n} e^{-  \beta |z|^{m}  }}{ S_{\alpha}(|z|^2)  }dA(z)    \\
    & =  \sum_{n= \frac{m}{2}  }^{+\infty}  
  \frac{ \alpha^{\frac{4n}{m}}}{\Gamma\left(\frac{2n+2}{m} \right)}   \left(\frac{t}{1+\alpha t}\right)^{\frac{2n+2}{m}}\int_{\mathbb C}  \frac{  |z|^{2n} e^{-  \beta |z|^m}  }{ S_{\alpha}(|z|^2)}   dA(z)   \\
  &  -  \sum_{n=\frac{m}{2} }^{+\infty}  
  \frac{ \beta^{\frac{4n}{m}}}{\Gamma\left(\frac{2n+2}{m} \right)}   \left(\frac{t}{ 1+\beta t}\right)^{\frac{2n+2}{m}}\int_{\mathbb C}  \frac{  |z|^{2n} e^{-  \alpha |z|^m}  }{ S_{\beta}(|z|^2)}    dA(z)
  \endaligned
  $$ 
  Dividing by $t^{\frac{2 +m}{m}}$ and taking the limit as $t$ tends to 0 in the latter equality yields
  $$  \aligned  &  \frac{ 2(\beta - \alpha ) }{m\Gamma\left(\frac{2}{m} \right)}  \int_{\mathbb C}  \frac{e^{-  \beta |z|^m}  }{ S_{\alpha}(|z|^2)}    dA(z)  \\
  &  =    
  \frac{ \alpha^{\frac{4n}{m}}}{\Gamma\left(\frac{2n+2}{m} \right)}  \int_{\mathbb C}  \frac{  |z|^{2n} e^{-  \beta |z|^m}  }{ S_{\alpha}(|z|^2)}   dA(z) 
  - \frac{ \beta^{\frac{4n}{m}}}{\Gamma\left(\frac{2n+2}{m} \right)}  \int_{\mathbb C}  \frac{  |z|^{2n} e^{-  \alpha |z|^m}  }{ S_{\beta}(|z|^2)}   dA(z).
    \endaligned
  $$ 
This shows that 
$$  \aligned 
& 
      \left(    
  \left(\frac{t}{1+\alpha t}\right)^{\frac{2}{m}} -  \left(\frac{t}{1+\beta t}\right)^{\frac{2}{m}}  -   \frac{2}{m} \left( \alpha - \beta \right) t^{\frac{m+2}{m}}   \right) 
  \frac{ 1}{\Gamma\left(\frac{2}{m} \right)}  \int_{\mathbb C}  
  \frac{  e^{-  \beta |z|^{m}  }}{ S_{\alpha}(|z|^2)  }dA(z) 
     \\
 & 
 + \sum_{n=1}^{\frac{m}{2} - 1}  
      \left(    
  \left(\frac{t}{1+\alpha t}\right)^{\frac{2n+2}{m}} -  \left(\frac{t}{1+\beta t}\right)^{\frac{m+2}{m}}  \right) 
  \frac{ \alpha^{\frac{4n}{m}}}{\Gamma\left(\frac{2n+2}{m} \right)}  \int_{\mathbb C}  
  \frac{  |z|^{2n} e^{-  \beta |z|^{m}  }}{ S_{\alpha}(|z|^2)  }dA(z) 
     \\
    + 
    & 
    \left( \left(\frac{t}{1+\alpha t}\right)^{\frac{m+2}{m}} -t^{\frac{m+2}{m}}\right) \frac{ \alpha^2}{\Gamma\left(\frac{m+2}{m} \right)}  \int_{\mathbb C}  \frac{  |z|^{2n} e^{-  \beta |z|^m}  }{ S_{\alpha}(|z|^2)}   dA(z)   \\
  & -
   \left( \left(\frac{t}{ 1+\beta t}\right)^{\frac{m+2}{m} } -t^{\frac{m+2}{m}}\right) \frac{ \beta^2}{\Gamma\left(\frac{m+2}{m} \right)}\int_{\mathbb C}  \frac{  |z|^{2n} e^{-  \alpha |z|^m}  }{ S_{\beta}(|z|^2)}    dA(z)
  \\
    & =  \sum_{n= \frac{m}{2} +1 }^{+\infty}  
  \frac{ \alpha^{\frac{4n}{m}}}{\Gamma\left(\frac{2n+2}{m} \right)}   \left(\frac{t}{1+\alpha t}\right)^{\frac{2n+2}{m}}\int_{\mathbb C}  \frac{  |z|^{2n} e^{-  \beta |z|^m}  }{ S_{\alpha}(|z|^2)}   dA(z)   \\
  &  -  \sum_{n=\frac{m}{2}  +1 }^{+\infty}  
  \frac{ \beta^{\frac{4n}{m}}}{\Gamma\left(\frac{2n+2}{m} \right)}   \left(\frac{t}{ 1+\beta t}\right)^{\frac{2n+2}{m}}\int_{\mathbb C}  \frac{  |z|^{2n} e^{-  \alpha |z|^m}  }{ S_{\beta}(|z|^2)}    dA(z)
  \endaligned
  $$ 
This shows that 
$$
\aligned
\int_{\mathbb C}  \frac{  |z|^{2n} e^{-  \beta |z|^m}  }{ S_{\alpha}(|z|^2)}  \int_{\mathbb C}   dA(z) = \int_{\mathbb C}  \frac{  |z|^{2n} e^{-  \alpha |z|^m}  }{ S_{\beta}(|z|^2)}    dA(z)
\endaligned
$$ 
for all integers $n\in [0, \frac{m}{2} +1[ \setminus\{\frac{m}{2}\}.$
An induction process allows us to establish the result of the lemma in the general case.
